\chapter{AST (Aspartate Aminotransferase)}

\section*{What Is This Test?}
Aspartate aminotransferase (AST), formerly called serum glutamic-oxaloacetic transaminase (SGOT), is an enzyme that catalyzes the reversible transfer of an amino group from aspartate to alpha-ketoglutarate, producing oxaloacetate and glutamate. AST is found in high concentrations in the liver (both cytoplasm and mitochondria), cardiac muscle, skeletal muscle, kidneys, brain, pancreas, lungs, red blood cells, and white blood cells. Because AST is widely distributed across multiple tissues, it is less specific for liver injury than ALT. In the hepatocyte, AST exists as two isoenzymes: cytoplasmic AST (approximately 20\% of total liver AST) and mitochondrial AST (approximately 80\%). In mild hepatocellular injury with plasma membrane damage, cytoplasmic AST and ALT are released, with ALT predominating. In more severe or chronic injury involving mitochondrial damage (such as alcoholic liver disease), mitochondrial AST is also released, elevating the AST:ALT ratio above 2:1. AST has a shorter serum half-life (approximately 17 hours) compared to ALT (approximately 47 hours). AST rises within hours of liver injury, peaks at 24--36 hours, and declines over 3--7 days. AST is measured as part of the comprehensive metabolic panel (CMP) and liver function tests (LFTs).

\section*{Alternative Names}
\begin{itemize}
  \item SGOT
  \item Serum Glutamic-Oxaloacetic Transaminase
  \item Aspartate Transaminase
  \item GOT
  \item Serum AST
\end{itemize}
\section*{Principle}
AST is measured by a kinetic enzymatic spectrophotometric method based on the reaction: L-aspartate + alpha-ketoglutarate $\rightarrow$ oxaloacetate + L-glutamate (catalyzed by AST). The oxaloacetate is reduced to malate by malate dehydrogenase (MDH) with concurrent oxidation of NADH to NAD\textsuperscript{+}: Oxaloacetate + NADH + H\textsuperscript{+} $\rightarrow$ Malate + NAD\textsuperscript{+}. The rate of decrease in absorbance at 340 nm is directly proportional to AST activity. This method was standardized by the IFCC. As with ALT, the assay requires addition of pyridoxal-5\textsuperscript{'}-phosphate (P5P) to ensure maximal enzyme activity. The test is performed at 37 degrees C on automated chemistry analyzers. Results are reported in IU/L. Critical pre-analytical note: hemolysis causes significant false elevation of AST because red blood cells contain AST at approximately 40 times the serum concentration (compared to approximately 5--6 times for ALT). Even slight hemolysis can substantially elevate AST. This is the most important pre-analytical consideration.

\section*{History}
AST was the first serum enzyme used in clinical diagnosis. Arthur Karmen, Felix Wr\'{o}blewski, and John LaDue demonstrated in 1954--1955 that serum AST was elevated in patients with acute myocardial infarction \cite{karmen1955transaminase}, opening the field of serum enzyme diagnostics. Before troponin, AST (along with CK and LDH) was a standard cardiac biomarker. The clinical utility of AST for liver disease was established by De Ritis in 1957 \cite{deritis1957enzymic}. The IFCC standardized method for AST was published alongside the ALT method.

\section*{Why This Test Is Ordered}
\begin{itemize}
  \item As part of the comprehensive metabolic panel for hepatocellular injury screening (always ordered with ALT)
  \item Evaluation of liver disease: acute viral hepatitis, alcoholic hepatitis (AST:ALT >2:1 is classic), drug-induced liver injury, ischemic hepatitis, NAFLD/NASH, autoimmune hepatitis
  \item Assessment of alcoholic liver disease: AST:ALT >2:1 strongly suggests alcohol as etiology. AST rarely exceeds 300--500 IU/L in alcoholic hepatitis (unlike viral hepatitis where AST can reach >1,000 IU/L)
  \item Diagnosis of ischemic hepatitis (shock liver): AST and ALT rise rapidly to >1,000 IU/L (often AST > ALT) in the setting of hypotension, cardiac arrest, sepsis; enzymes fall rapidly with restoration of perfusion
  \item Detection of muscle injury: AST is elevated in rhabdomyolysis, polymyositis/dermatomyositis, muscular dystrophies, and after vigorous exercise. Creatine kinase (CK) is the preferred test for muscle injury; AST alone is non-specific
  \item Adjunctive marker in acetaminophen toxicity and other drug-induced liver injuries
  \item Evaluation of Wilson disease: AST:ALT ratio >2:1 is seen in approximately 80\% of Wilson disease patients
\end{itemize}

\section*{Primary vs Secondary Test}
AST is a primary screening test as part of the CMP, always interpreted alongside ALT. AST alone is non-specific; the AST:ALT ratio provides etiologic clues.

\section*{How This Test Is Performed}
A healthcare professional draws blood from a vein into a serum separator tube (gold-top) or lithium heparin tube (green-top). Hemolysis MUST be avoided---even slightly hemolyzed samples will produce falsely elevated AST and should be rejected. The tourniquet should be released promptly to minimize stasis and hemolysis. The sample is centrifuged and AST measured on an automated chemistry analyzer. Results are available within 30--60 minutes.

\section*{What Preparation Is Needed}
Same considerations as ALT: 8--12 hour fasting recommended when part of CMP. Alcohol intake within 24--48 hours elevates AST. Hepatotoxic medications (acetaminophen, isoniazid, valproic acid) and herbal supplements (green tea extract, kava) can elevate AST. Strenuous exercise within 24--48 hours can cause mild AST elevation from muscle release (more pronounced for AST than ALT due to greater muscle AST content). Intramuscular injections can cause mild AST release.

\section*{Contraindications and Precautions}
\begin{itemize}
  \item No absolute contraindications. Critical pre-analytical considerations:
  \item (1) Hemolysis is the single most important interferent---AST concentration in erythrocytes is approximately 40$\times$ the serum level (compared to 5--6$\times$ for ALT). Even slight hemolysis increases AST. Hemolyzed specimens must be rejected.
  \item (2) Pyridoxine (vitamin B6) deficiency: AST, like ALT, requires P5P as a cofactor. Alcoholics and malnourished patients may have falsely low AST due to B6 deficiency unless the assay reagents include excess P5P.
  \item (3) Extensive muscle injury (rhabdomyolysis, surgery, IM injections, prolonged seizures, crush trauma) elevates AST substantially, masking true liver-derived AST changes. Concurrent CK measurement distinguishes muscle from liver source.
  \item (4) Macro-AST: a rare phenomenon where AST forms a complex with an immunoglobulin (usually IgG), resulting in persistently elevated AST because the complex is not cleared from circulation. Suggested by isolated AST elevation without ALT elevation and no clinical liver or muscle disease. Polyethylene glycol (PEG) precipitation testing confirms.
  \item (5) AST may be mildly elevated in pregnancy (third trimester) from placental release.
  \item (6) AST declines slightly with age.
\end{itemize}
\section*{Risks}
Minimal risk from venipuncture. Mild pain or bruising at site resolves in 1--2 days. Rare: hematoma, infection, nerve injury, vasovagal syncope.

\section*{What Happens After the Test}
Results available within 30--60 minutes. AST is always interpreted simultaneously with ALT via the De Ritis ratio (AST:ALT): AST:ALT >2:1: strongly suggests alcoholic liver disease (alcohol damages mitochondria, releasing mitochondrial AST; alcohol also depletes pyridoxine, which preferentially affects ALT synthesis). AST:ALT >2:1 is also seen in Wilson disease (80\% of cases), and in cirrhosis of any cause (reduced ALT synthesis). AST:ALT <1: typical of viral hepatitis, NAFLD, and drug-induced liver injury. AST:ALT approaching 1:1 with rising bilirubin may indicate progression to cirrhosis. AST >10,000 IU/L: typically acetaminophen toxicity or severe ischemic hepatitis. AST >3,000 IU/L with ALT proportionally elevated: acute viral hepatitis or autoimmune hepatitis. Without ALT elevation, an isolated AST elevation suggests muscle source, macro-AST, or hemolysis artifact.

\section*{Understanding Results}

\subsection*{Below Normal}
\begin{itemize}
  \item Low AST is not clinically significant and does not indicate organ dysfunction.
  \item Uremia (CKD/ESRD) may suppress AST activity in vitro, producing falsely low values.
  \item Pyridoxine (vitamin B6) deficiency in alcoholics, malnourished patients, or elderly reduces hepatic AST synthesis.
  \item Pregnancy (third trimester): mild decrease from hemodilution.
  \item In end-stage liver disease, AST may paradoxically be low despite cirrhosis.
\end{itemize}

\subsection*{Above Normal}
\begin{itemize}
  \item \textbf{Isolated AST elevation (ALT normal):} Suspect muscle source (rhabdomyolysis, intense exercise, IM injections, myopathy, seizure, trauma), macro-AST, hemolysis artifact, or early alcoholic liver disease (AST rises before ALT in early alcohol injury).
  \item \textbf{Alcoholic liver disease (AST:ALT >2:1):} AST rarely exceeds 300--500 IU/L even in severe alcoholic hepatitis---an AST >500 IU/L with AST:ALT >2:1 suggests a co-existing insult (acetaminophen, viral hepatitis, ischemic).
  \item \textbf{Acute viral hepatitis (HAV, HBV):} AST and ALT rise to 1,000--5,000 IU/L with AST:ALT <1 (ALT predominates). AST peaks at 24--36 hours, ALT peaks at 1--3 days. ALT remains elevated longer than AST.
  \item \textbf{Acetaminophen toxicity:} Marked AST/ALT elevation (3,000--10,000+ IU/L at 24--72 hours post-ingestion). AST and ALT are similarly elevated. AST >10,000 IU/L with ALT >5,000 IU/L is classic. Hyperacute transaminitis with very rapid rise.
  \item \textbf{Ischemic hepatitis (shock liver):} Rapid rise of AST and ALT (>1,000 IU/L, AST often > ALT) within 24 hours of hypotensive event (cardiac arrest, sepsis, severe hemorrhage). Enzymes normalize rapidly (within 3--7 days) once perfusion is restored. LDH is also markedly elevated. This is one of the most common causes of massive transaminase elevation in hospitalized patients.
  \item \textbf{Wilson disease:} AST:ALT >2:1 is seen in approximately 80\% of cases. Patients typically <40 years with neurological symptoms, psychiatric symptoms, or unexplained liver disease. Serum ceruloplasmin <20 mg/dL, 24-hour urine copper >100 mcg/day.
  \item \textbf{Muscle disorders:} Rhabdomyolysis (CK >5--10$\times$ ULN, AST elevated, ALT may be normal or mildly elevated), polymyositis/dermatomyositis (CK elevated), muscular dystrophy (CK markedly elevated from birth), recent vigorous exercise (transient, resolves in 48--72 hours), intramuscular injections, seizures, crush injury.
  \item \textbf{Other causes:} Celiac disease (reversible transaminitis in up to 40\% at diagnosis), hemochromatosis, alpha-1 antitrypsin deficiency, autoimmune hepatitis, HELLP syndrome (hemolysis, elevated liver enzymes, low platelets in pregnancy/preeclampsia), heat stroke.
\end{itemize}

\section*{Normal Results}
\begin{itemize}
  \item \textbf{Adult males:} 10--40 IU/L
  \item \textbf{Adult females:} 9--32 IU/L
  \item \textbf{Children:} 10--50 IU/L (age-dependent)
  \item \textbf{Infants:} 15--60 IU/L
  \item \textbf{Newborns:} 25--75 IU/L (higher at birth)
  \item \textbf{AST half-life:} 17 hours (shorter than ALT's 47 hours)
  \item \textbf{AST:ALT ratio (De Ritis ratio):} Normally approximately 0.8:1
\end{itemize}

\section*{Follow-up Tests}
\begin{itemize}
  \item \textbf{ALT:} Always ordered together. AST:ALT ratio guides differential diagnosis.
  \item \textbf{ALP, GGT, total and direct bilirubin, albumin, INR:} Complete liver panel for synthetic function and cholestasis assessment.
  \item \textbf{Creatine kinase (CK):} If muscle source is suspected as cause of AST elevation. CK >5$\times$ ULN with normal ALT confirms muscle origin.
  \item \textbf{Viral hepatitis serologies:} HAV IgM, HBsAg, anti-HBc IgM, anti-HCV.
  \item \textbf{Abdominal ultrasound with Doppler:} First-line imaging for liver disease.
  \item \textbf{Ceruloplasmin, 24-hour urine copper:} Wilson disease workup if AST:ALT >2 and patient <40 years.
  \item \textbf{Acetaminophen level:} If toxicity suspected.
  \item \textbf{Serum iron studies:} Ferritin, transferrin saturation.
  \item \textbf{Liver biopsy/fibroscan:} For fibrosis staging.
\end{itemize}

\section*{References}
\bibliographystyle{unsrtnat}
\bibliography{references}

\clearpage
\section*{Regulatory Status and Authorization Date}
This chapter describes a test category, not a specific commercial product. FDA, CDSCO, EU IVDR, and MHRA approval or registration dates therefore cannot be assigned without the assay manufacturer, product name, instrument, and intended use. For laboratory-developed tests, an individual agency approval date may not exist. See the Preface for the interpretation of regulatory status.

\clearpage
